\documentclass[11pt]{article}
\usepackage[margin=1in]{geometry}
\usepackage{amsmath,amssymb,amsthm,mathtools}
\usepackage{bm}
\usepackage{hyperref}

\title{A Pure Probability Formulation of the Centered Gaussian Bridge Problem}
\author{}
\date{\today}

\newcommand{\R}{\mathbb{R}}
\newcommand{\Pbb}{\mathbb{P}}
\newcommand{\N}{\mathcal{N}}
\newcommand{\ip}[2]{\left\langle #1,#2\right\rangle}
\newcommand{\norm}[1]{\left\lVert #1\right\rVert}
\newcommand{\PhiStd}{\Phi}
\newcommand{\Beps}{\mathcal{B}_\varepsilon}

\newtheorem{theorem}{Theorem}
\newtheorem{lemma}{Lemma}
\newtheorem{corollary}{Corollary}

\begin{document}
\maketitle

\section{The object under study}
Fix dimension $n\ge 1$, noise scale $\sigma>0$, and a vector $d\in\R^n$.
Consider two Gaussian laws on $\R^n$:
\[
Y\sim \N(d,\sigma^2 I_n),
\qquad
Y'\sim \N(-d,\sigma^2 I_n).
\]
These are two equally spread clouds with opposite centers.

Let $p_d,p_{-d}$ be their densities. For measurable $S\subseteq\R^n$, define
\[
\Delta_\varepsilon(S)
:=\Pbb(Y\in S)-e^{\varepsilon}\Pbb(Y'\in S),
\qquad \varepsilon\in\R.
\]
The central quantity is
\[
\sup_{S\in\mathcal B(\R^n)} \Delta_\varepsilon(S).
\]
This is a weighted set-separation functional between two Gaussian laws.

\section{Equivalent statistical views of the same quantity}

\subsection*{View A: extremal measurable set}
A standard variational identity gives
\[
\sup_S \Delta_\varepsilon(S)
=\int_{\R^n}(p_d-e^{\varepsilon}p_{-d})_+.
\]
Hence the maximizing set is
\[
S_\varepsilon^*:=\{y\in\R^n: p_d(y)-e^{\varepsilon}p_{-d}(y)\ge 0\}.
\]
So the problem is exactly: characterize this set and evaluate the corresponding integral.

\subsection*{View B: log-likelihood geometry}
The log-likelihood ratio is
\[
\log\frac{p_d(y)}{p_{-d}(y)}=\frac{2\ip{d}{y}}{\sigma^2}.
\]
Therefore
\[
S_\varepsilon^*
=\left\{y:\ip{d}{y}\ge \frac{\varepsilon\sigma^2}{2}\right\},
\]
a halfspace orthogonal to $d$. Intuitively: all nonlinear coordinates are irrelevant; only projection onto $d$ matters.

\subsection*{View C: one-dimensional reduction}
Define the scalar projection
\[
T(y):=-\ip{d}{y}.
\]
Then
\[
T(Y)\sim \N(-r,\sigma^2 r),
\qquad
T(Y')\sim \N(+r,\sigma^2 r),
\qquad
r:=\norm{d}^2.
\]
And the extremal halfspace becomes a threshold event:
\[
\ip{d}{y}\ge \frac{\varepsilon\sigma^2}{2}
\iff
T(y)\le -\frac{\varepsilon\sigma^2}{2}.
\]
So the original $n$-dimensional supremum is equivalent to two one-dimensional Gaussian CDF values.

\section{What must be matched exactly}
Set
\[
\tau:=-\frac{\varepsilon\sigma^2}{2},
\qquad r=\norm{d}^2.
\]
Then (for $r>0$)
\[
\Pbb(T(Y)\le\tau)=
\PhiStd\!\left(\frac{\tau+r}{\sigma\sqrt r}\right),
\qquad
\Pbb(T(Y')\le\tau)=
\PhiStd\!\left(\frac{\tau-r}{\sigma\sqrt r}\right).
\]
After simplification,
\[
\frac{\tau+r}{\sigma\sqrt r}=\frac{\sqrt r}{\sigma}-\frac{\varepsilon\sigma}{2\sqrt r},
\qquad
\frac{\tau-r}{\sigma\sqrt r}=-\frac{\sqrt r}{\sigma}-\frac{\varepsilon\sigma}{2\sqrt r}.
\]
If one introduces
\[
\Delta:=2\sqrt r,
\]
these become
\[
\frac{\Delta}{2\sigma}-\frac{\varepsilon\sigma}{\Delta},
\qquad
-\frac{\Delta}{2\sigma}-\frac{\varepsilon\sigma}{\Delta}.
\]
So the same statistical quantity can be parameterized either by $r=\norm d^2$ or by $\Delta=2\norm d$.

\section{Exact bridge endpoints (tightened)}
Define the \emph{Gaussian bridge functional}
\[
\Beps(d;\sigma)
:=
\sup_{S\in\mathcal B(\R^n)}
\left[\Pbb(Y\in S)-e^\varepsilon \Pbb(Y'\in S)\right],
\]
for $Y\sim\N(d,\sigma^2I_n)$ and $Y'\sim\N(-d,\sigma^2I_n)$.

Define the scalar profile (for $r>0$)
\[
G_{\varepsilon,\sigma}(r)
:=
\PhiStd\!\left(\frac{\sqrt r}{\sigma}-\frac{\varepsilon\sigma}{2\sqrt r}\right)
-e^\varepsilon
\PhiStd\!\left(-\frac{\sqrt r}{\sigma}-\frac{\varepsilon\sigma}{2\sqrt r}\right),
\]
and $G_{\varepsilon,\sigma}(0):=(1-e^\varepsilon)_+$.

Define also the $\Delta$-parameterized profile
\[
H_{\varepsilon,\sigma}(\Delta)
:=
\PhiStd\!\left(\frac{\Delta}{2\sigma}-\frac{\varepsilon\sigma}{\Delta}\right)
-e^\varepsilon
\PhiStd\!\left(-\frac{\Delta}{2\sigma}-\frac{\varepsilon\sigma}{\Delta}\right),
\qquad \Delta>0.
\]

Define two endpoint propositions:
\begin{itemize}
\item \textbf{(A) Radius-constrained vector statement}
\[
\forall d\in\R^n,\ \norm d\le \frac{\Delta}{2}
\implies
\Beps(d;\sigma)\le \delta.
\]
\item \textbf{(B) Single-scalar endpoint statement}
\[
H_{\varepsilon,\sigma}(\Delta)\le \delta.
\]
\end{itemize}

These are deliberately adjacent: both are now inequalities about one explicit
functional family, only indexed differently (vector radius versus scalar $\Delta$).

\section{Minimal equivalence to prove}
Using standard Gaussian likelihood-ratio calculus and one-dimensional projection,
the bridge reduces to two analysis statements:
\begin{enumerate}
\item \textbf{Radial identity (rotation reduction):}
\[
\Beps(d;\sigma)=G_{\varepsilon,\sigma}(\norm d^2)\quad\text{for all }d\in\R^n.
\]
\item \textbf{Endpoint envelope identity:}
\[
\sup_{0\le r\le (\Delta/2)^2} G_{\varepsilon,\sigma}(r)
=
H_{\varepsilon,\sigma}(\Delta).
\]
\end{enumerate}
Since $H_{\varepsilon,\sigma}(\Delta)=G_{\varepsilon,\sigma}(\Delta^2/4)$ for
$\Delta>0$, the second line is exactly a one-variable envelope statement.
In other words, once radial reduction is accepted, the remaining question is purely:
\[
\text{how does }G_{\varepsilon,\sigma}(r)\text{ behave on }[0,(\Delta/2)^2]?
\]

\section{Sharp bridge seam}
For a standard PhD-level probability reader, the seam is exactly:
\[
\text{(A)}
\Longleftrightarrow
\Big(\sup_{0\le r\le(\Delta/2)^2} G_{\varepsilon,\sigma}(r)\le \delta\Big)
\Longleftrightarrow
\text{(B)}.
\]
Hence the remaining gap is one scalar fact:
\[
\sup_{0\le r\le R} G_{\varepsilon,\sigma}(r)=G_{\varepsilon,\sigma}(R),
\qquad R=(\Delta/2)^2,
\]
i.e.\ monotonicity of one explicit one-variable profile.

\section{Current formal seam (Lean-oriented)}
The actual formal target is the theorem-level implication in
\texttt{Mf/DP/GaussianCore.lean} that upgrades an explicit scalar bound into the
set-wise Gaussian bound for all measurable events. Concretely, the bridge is:
\begin{itemize}
\item \textbf{End A (scalar endpoint):} an explicit inequality at $\Delta$
\[
\PhiStd\!\left(\frac{\Delta}{2\sigma}-\frac{\varepsilon\sigma}{\Delta}\right)
-e^\varepsilon
\PhiStd\!\left(-\frac{\Delta}{2\sigma}-\frac{\varepsilon\sigma}{\Delta}\right)
\le \delta;
\]
\item \textbf{End B (event endpoint):} for all measurable $S\subseteq\R^n$,
\[
\Pbb(Y\in S)-e^\varepsilon \Pbb(Y'\in S)\le \delta
\]
for the centered Gaussian pair $Y\sim\N(d,\sigma^2I_n)$,
$Y'\sim\N(-d,\sigma^2I_n)$ under $\|d\|\le \Delta/2$.
\end{itemize}

For the formal development, these are used under the natural domain assumptions
\[
\sigma>0,\qquad \Delta\ge 0.
\]
These are not cosmetic: they are exactly what makes the reparameterization and
the comparison map
\[
\Delta \longmapsto \frac{\Delta}{2\sigma}
\]
well-behaved and order-preserving.

The remaining formal step is therefore best stated as:
\begin{itemize}
\item from End A, i.e.
\[
\PhiStd\!\left(\frac{\Delta}{2\sigma}-\frac{\varepsilon\sigma}{\Delta}\right)
-e^\varepsilon
\PhiStd\!\left(-\frac{\Delta}{2\sigma}-\frac{\varepsilon\sigma}{\Delta}\right)
\le \delta,
\]
\item and from
\[
r:=\norm d^2,\qquad r\le (\Delta/2)^2,
\]
\item deduce the same bound at the exact radius parameter
\[
\Delta_r:=2\sqrt r
\]
by monotonicity of the scalar profile on $(0,\infty)$, then transfer to End B.
\end{itemize}

Then one closes by a clean case split:
\begin{itemize}
\item $r>0$: use the explicit two-$\sqrt r$ reduction;
\item $r=0$: the two Gaussian laws coincide, so the bridge value is the trivial
degenerate edge case.
\end{itemize}

So the final gap is no longer conceptual; it is a precise endpoint transfer
lemma under $\sigma>0,\Delta\ge 0$, followed by the $r=0$/$r>0$ split.

\section{Monotonicity of the scalar profile \texorpdfstring{$G_{\varepsilon,\sigma}$}{G}}
Set
\[
a:=\frac{\sqrt r}{\sigma}>0,
\quad
u(a):=a-\frac{\varepsilon}{2a},
\quad
v(a):=-a-\frac{\varepsilon}{2a},
\]
so that for $r>0$,
\[
G_{\varepsilon,\sigma}(r)
=\PhiStd\!\bigl(u(a)\bigr)-e^\varepsilon\PhiStd\!\bigl(v(a)\bigr).
\]

\begin{lemma}[Gaussian density symmetry]\label{lem:density-sym}
For all $a>0$,
\[
e^\varepsilon\,\varphi\!\bigl(v(a)\bigr)=\varphi\!\bigl(u(a)\bigr),
\qquad
\varphi(x):=\frac{1}{\sqrt{2\pi}}e^{-x^2/2}.
\]
\end{lemma}

\begin{proof}
A direct computation shows $v(a)^2-u(a)^2=2\varepsilon$, hence
\[
\varphi\!\bigl(v(a)\bigr)=\varphi\!\bigl(u(a)\bigr)e^{-\varepsilon}.
\]
Multiplying by $e^\varepsilon$ yields the claim.
\end{proof}

\begin{theorem}[Strict monotonicity]\label{thm:mono}
For every fixed $\varepsilon\in\R$ and $\sigma>0$ the map
$r\mapsto G_{\varepsilon,\sigma}(r)$ is
\begin{enumerate}
\item continuous on $[0,\infty)$,
\item strictly increasing on $(0,\infty)$.
\end{enumerate}
Consequently,
\[
\sup_{0\le r\le (\Delta/2)^2}G_{\varepsilon,\sigma}(r)
=
G_{\varepsilon,\sigma}\!\bigl((\Delta/2)^2\bigr)
\quad (\Delta\ge 0),
\]
and for $\Delta>0$,
\[
H_{\varepsilon,\sigma}(\Delta)=G_{\varepsilon,\sigma}(\Delta^2/4).
\]
\end{theorem}

\begin{proof}
Write $F(a):=\PhiStd(u(a))-e^\varepsilon\PhiStd(v(a))$ so that
$G_{\varepsilon,\sigma}(r)=F(\sqrt r/\sigma)$ for $r>0$.
Differentiate:
\[
F'(a)=\varphi(u(a))u'(a)-e^\varepsilon\varphi(v(a))v'(a).
\]
Using Lemma~\ref{lem:density-sym} and
\[
u'(a)=1+\frac{\varepsilon}{2a^2},
\qquad
v'(a)=-1+\frac{\varepsilon}{2a^2},
\]
gives
\[
F'(a)=\varphi(u(a))\bigl(u'(a)-v'(a)\bigr)
=2\,\varphi(u(a))>0
\quad(a>0),
\]
so $F$, hence $G_{\varepsilon,\sigma}$, is strictly increasing on $(0,\infty)$.

For continuity at $r=0$, let $a\downarrow0$:
\begin{itemize}
\item if $\varepsilon>0$ then $u(a),v(a)\to-\infty$ and $F(a)\to0$;
\item if $\varepsilon<0$ then $u(a),v(a)\to+\infty$ and $F(a)\to 1-e^\varepsilon$;
\item if $\varepsilon=0$ then $F(a)=2\PhiStd(a)-1\to0$.
\end{itemize}
Thus $\lim_{r\downarrow0}G_{\varepsilon,\sigma}(r)=(1-e^\varepsilon)_+$,
the declared value at $r=0$.
\end{proof}

\begin{corollary}[Envelope identity]\label{cor:envelope}
For every $\Delta\ge0$,
\[
\sup_{0\le r\le (\Delta/2)^2} G_{\varepsilon,\sigma}(r)
=
\begin{cases}
H_{\varepsilon,\sigma}(\Delta), & \Delta>0,\\[6pt]
(1-e^\varepsilon)_+, & \Delta=0.
\end{cases}
\]
\end{corollary}

\begin{proof}
Use Theorem~\ref{thm:mono}: the supremum on a closed interval of a strictly
increasing function is attained at the right endpoint.
\end{proof}

\section{Closing the bridge}

\begin{theorem}[Endpoint transfer]\label{thm:bridge}
Assume $\sigma>0$ and $\Delta\ge0$.
The radius-constrained statement \textbf{(A)} is equivalent to the
single-scalar statement \textbf{(B)}:
\[
\forall d,\;\|d\|\le\frac{\Delta}{2}\;\Longrightarrow\;
\Beps(d;\sigma)\le\delta
\quad\Longleftrightarrow\quad
H_{\varepsilon,\sigma}(\Delta)\le\delta.
\]
\end{theorem}

\begin{proof}
($\Rightarrow$)
Take $d$ with $\|d\|=\Delta/2$ (possible if $\Delta>0$; otherwise $d=0$).
Apply \textbf{(A)} and the radial identity
\[
\Beps(d;\sigma)=G_{\varepsilon,\sigma}(\|d\|^2)
\]
to get
\[
G_{\varepsilon,\sigma}(\Delta^2/4)\le\delta,
\]
i.e.\ $H_{\varepsilon,\sigma}(\Delta)\le\delta$.

($\Leftarrow$)
Assume $H_{\varepsilon,\sigma}(\Delta)\le\delta$.
For any $d$ with $\|d\|\le\Delta/2$,
\[
\Beps(d;\sigma)=G_{\varepsilon,\sigma}(\|d\|^2)
\le
\sup_{0\le r\le(\Delta/2)^2}G_{\varepsilon,\sigma}(r)
=H_{\varepsilon,\sigma}(\Delta)
\le\delta,
\]
using Corollary~\ref{cor:envelope}.
\end{proof}

\section{What this solves analytically}
With Theorem~\ref{thm:bridge}, the analytic seam is closed: verifying a
single scalar inequality at $\Delta$ is both necessary and sufficient for the
bound to hold for every vector $d$ with $\|d\|\le\Delta/2$.

\medskip\noindent
\textbf{Practical upshot.}
All the high-dimensional, set-valued optimization collapses to checking
\[
H_{\varepsilon,\sigma}(\Delta)\le\delta,
\]
a one-line condition involving just the standard normal CDF.

\end{document}
